\documentclass[11pt,3p,review]{elsarticle}
\usepackage[utf8]{inputenc}
\usepackage{amsmath,amssymb,amsfonts}
\usepackage{booktabs}
\usepackage{graphicx}
\usepackage{hyperref}

\journal{Artificial Intelligence}

\begin{document}

\begin{frontmatter}

\title{StateShift: Tracking State-Dependent Reasoning Recovery Across Post-Training}

\author[inst1]{Author Name\corref{cor1}}
\ead{REQUIRED_FROM_AUTHOR@institution.edu}
\cortext[cor1]{Corresponding author. Full postal address required.}
\address[inst1]{Department Name, Institution Name, City, Country}

\begin{abstract}
Standard evaluations of reinforcement-learning post-training in large language models emphasize aggregate benchmark accuracy, potentially obscuring changes conditional on local reasoning state. We introduce \textbf{StateShift}, a framework for measuring target-transition recovery under matched reasoning states. In a controlled study ($N=454, 29,056$ rollouts), we observe an 11.76-percentage-point state-by-checkpoint interaction between the base and step-256 checkpoints ($\Gamma_{256}=0.1176, 95\%$ problem-blocked bootstrap CI $[0.0955, 0.1400]$). Across nine empirically evaluated checkpoints, the interaction was consistent with a non-decreasing trajectory under prespecified order-restricted analysis despite local variation in unconstrained estimates, and was already detectable at the earliest available post-training checkpoint, $t=32$ ($\Gamma_{32}=0.0333$, multiplicity-adjusted 95\% CI $[0.0011, 0.0655]$). In a separate analysis of 3,200 unperturbed rollouts, 18.19\% contained a verifier-confirmed natural reasoning error; among 582 qualifying error episodes, 180 subsequently satisfied the prespecified autonomous recovery criterion (30.93\%, 95\% problem-blocked bootstrap CI $[27.19\%, 34.82\%]$). These results show that state-conditioned evaluation exposes behavioral structure not captured by aggregate correctness alone.
\end{abstract}

\begin{keyword}
Large language models \sep Reinforcement learning \sep Mathematical reasoning \sep Reasoning recovery \sep Post-training evaluation \sep State-conditioned evaluation
\end{keyword}

\end{frontmatter}

\section{Introduction}
Reinforcement learning from verifier feedback (RLVR) has significantly improved mathematical reasoning in large language models. However, standard benchmark metrics measure only final accuracy, providing limited insight into how post-training changes reasoning behavior conditional on local state. Does RL training uniformly improve reasoning capability, or does it specifically enhance the model's ability to recover from invalid intermediate states?

To answer this question, we present \textbf{StateShift}, a experimental framework evaluating state-conditioned reasoning dynamics across post-training.

\section{Background and Related Work}
\subsection{Reinforcement Learning for Reasoning}
Recent advances such as DeepSeekMath, DeepSeek-R1, and DeepScaler demonstrate the efficacy of Group Relative Policy Optimization (GRPO) for reasoning.

\subsection{Categorical Self-Correction Literature}
Self-correction literature generally requires multi-turn verbal feedback or PRM verifiers. In contrast, StateShift evaluates unprompted natural post-error recovery in single-rollout decoding.

\section{The StateShift Framework}
StateShift evaluates target-transition success under two matched conditions:
\begin{enumerate}
    \item \textbf{Recovery Condition ($R$)}: The model is initialized at a locally invalid intermediate reasoning state.
    \item \textbf{Control Condition ($C$)}: The model is initialized at a matched control state.
\end{enumerate}

The primary estimand is the difference-in-differences state-by-checkpoint interaction:
\begin{equation}
\Gamma_t = (\mu_{R,t} - \mu_{R,0}) - (\mu_{C,t} - \mu_{C,0})
\end{equation}

\section{Controlled State-Conditioned Evaluation (Study A)}
In our primary confirmatory experiment ($N=454, K=16, 29,056$ rollouts), we observe:
\begin{align}
\mu_{R,0} = 0.3834, \quad \mu_{C,0} = 0.3892 \\
\mu_{R,256} = 0.7039, \quad \mu_{C,256} = 0.5921
\end{align}
yielding a primary interaction $\Gamma_{256} = \mathbf{+0.1176}$ ($95\%$ problem-blocked bootstrap CI $[+0.0955, +0.1400], p < 0.0001$). Under strict decontamination filtering ($N_{\text{Strict}}=388$), $\Gamma_{256,\text{Strict}} = \mathbf{+0.1160}$ ($95\%$ CI $[+0.0913, +0.1408]$).

\section{Post-Training Trajectory Analysis}
Evaluating intermediate checkpoints across training yields the empirical nine-point vector:
\begin{equation}
\mathbf{\Gamma} = [0.0000, +0.0333, +0.0337, +0.0774, +0.0748, +0.0598, +0.0976, +0.0950, +0.1176]
\end{equation}
Across nine empirically evaluated checkpoints, the interaction was consistent with a non-decreasing trajectory under prespecified order-restricted analysis despite local variation in unconstrained estimates. The positive interaction was already statistically detectable by the earliest available post-training checkpoint $t=32$ ($\Gamma_{32} = +0.0333$, multiplicity-adjusted 95\% CI $[+0.0011, +0.0655]$).

\section{Natural Post-Error Recovery (Study B)}
In $3,200$ unperturbed rollouts ($N=200, K=16$), $18.19\%$ contained a verifier-confirmed natural reasoning error ($\text{NEI} = 18.19\%$, $582$ error rollouts). Among $582$ qualifying error episodes, $180$ subsequently satisfied the prespecified autonomous recovery criterion ($\text{NRR} = \mathbf{30.93\%}$, $95\%$ problem-blocked bootstrap CI $[27.19\%, 34.82\%]$).

\section{Limitations}
Strict pairwise monotonicity across all adjacent intervals was not established due to sample size constraints ($K=2/3$ at intermediate steps). Sub-32-step emergence remains unobserved because fine-grained checkpoints $t \in \{1..31\}$ do not exist. Natural post-error recovery is conditional on error occurrence and should not be interpreted as an unqualified causal self-correction probability.

\section{Declarations}
\subsection{Competing Interests}
The authors declare that they have no known competing financial interests or personal relationships that could have appeared to influence the work reported in this paper.

\subsection{Funding}
No external funding supported this research.

\subsection{CRediT Author Statement}
Author Name: Conceptualization, Methodology, Software, Formal analysis, Writing - Original Draft.

\subsection{Data Availability}
The raw rollout records, analysis scripts, and problem registries supporting the findings of this study are available in the project repository.

\bibliographystyle{elsarticle-num}
\bibliography{03_References}

\end{document}
